\begin{center}
\textbf{Preços agrícolas, capacidade produtiva e perdas anteriores na conversão da vegetação nativa em Mato Grosso}

\vspace{0.4cm}
\textbf{Agricultural prices, productive capacity and prior crop losses in native vegetation conversion in Mato Grosso}
\end{center}

\vspace{0.5cm}
\noindent\textbf{Resumo}

\vspace{0.3cm}
A expansão agropecuária pode aproveitar áreas abertas ou converter vegetação, com implicações distintas para a gestão ambiental. Este artigo examina a associação dos preços reais de soja e milho com a conversão bruta para agricultura em Mato Grosso, considerando a capacidade produtiva potencial e as perdas anteriores. O painel municipal de 2001--2024 combina dados ambientais, produtivos, de preços e clima, com efeitos fixos municipais e de bioma por ano. Os gradientes preço×capacidade diminuem quando se exige histórico suficiente para medir perdas. Para aumento de 10\% no preço e dez pontos percentuais adicionais de perda, os contrastes são 0,0040 p.p. na soja e $-0,0001$ p.p. no milho, com intervalos de 95\% que abrangem zero nas três inferências. Comparações entre períodos e destinos revelam heterogeneidade sensível à amostra e à inferência. As exposições territoriais não identificam efeitos do Código Florestal ou da Moratória da Soja. Os resultados ressaltam a necessidade de distinguir destinos da abertura, condições produtivas e cobertura dos dados na avaliação das pressões sobre a vegetação.

\vspace{0.4cm}
\noindent\textbf{Palavras-chave:} conversão da vegetação, preços agrícolas, capacidade produtiva, quebras de safra, gestão ambiental.

\vspace{0.6cm}
\noindent\textbf{Abstract}

\vspace{0.3cm}
\begin{otherlanguage}{english}
Agricultural expansion can use cleared land or convert native vegetation, with different implications for environmental management. This article examines the association between real soybean and maize prices and gross conversion to agriculture in Mato Grosso, considering potential productive capacity and prior crop losses. The 2001--2024 municipal panel combines environmental, production, price and climate data, with municipality and biome-by-year fixed effects. Price--capacity gradients decrease when sufficient historical data are required to measure losses. For a 10\% price increase and ten additional percentage points of loss, the contrasts are 0.0040 percentage points for soybeans and $-0.0001$ for maize, with 95\% intervals spanning zero under all three inference procedures. Comparisons across periods and land-use destinations reveal heterogeneity sensitive to the sample and inference procedure. Territorial exposures do not identify effects of the Forest Code or the Soy Moratorium. The results highlight the need to distinguish conversion destinations, productive conditions and data coverage when assessing pressures on native vegetation.

\vspace{0.4cm}
\noindent\textbf{Keywords:} vegetation conversion, agricultural prices, productive capacity, crop losses, environmental management.
\end{otherlanguage}

\vspace{0.4cm}
\noindent\textbf{Classificação JEL:} Q11; Q15; Q57; C23; R14.
